\documentclass{amsart}
% For dealing with references we use the comment environment
\usepackage{verbatim}
\newenvironment{reference}{\comment}{\endcomment}
%\newenvironment{reference}{}{}
\newenvironment{slogan}{\comment}{\endcomment}
\newenvironment{history}{\comment}{\endcomment}

% For commutative diagrams we use Xy-pic
\usepackage[all]{xy}

% We use 2cell for 2-commutative diagrams.
\xyoption{2cell}
\UseAllTwocells

% We use multicol for the list of chapters between chapters
\usepackage{multicol}

% This is generally recommended for better output
\usepackage{lmodern}
\usepackage[T1]{fontenc}

% For cross-file-references

% Package for hypertext links:
\usepackage{hyperref}

% For any local file, say "hello.tex" you want to link to please

% Theorem environments.
%
\theoremstyle{plain}
\newtheorem{theorem}[subsection]{Theorem}
\newtheorem{proposition}[subsection]{Proposition}
\newtheorem{lemma}[subsection]{Lemma}

\theoremstyle{definition}
\newtheorem{definition}[subsection]{Definition}
\newtheorem{example}[subsection]{Example}
\newtheorem{exercise}[subsection]{Exercise}
\newtheorem{situation}[subsection]{Situation}

\theoremstyle{remark}
\newtheorem{remark}[subsection]{Remark}
\newtheorem{remarks}[subsection]{Remarks}

\numberwithin{equation}{subsection}

% Macros
%
\def\lim{\mathop{\mathrm{lim}}\nolimits}
\def\colim{\mathop{\mathrm{colim}}\nolimits}
\def\Spec{\mathop{\mathrm{Spec}}}
\def\Hom{\mathop{\mathrm{Hom}}\nolimits}
\def\Ext{\mathop{\mathrm{Ext}}\nolimits}
\def\SheafHom{\mathop{\mathcal{H}\!\mathit{om}}\nolimits}
\def\SheafExt{\mathop{\mathcal{E}\!\mathit{xt}}\nolimits}
\def\Sch{\mathit{Sch}}
\def\Mor{\mathop{\mathrm{Mor}}\nolimits}
\def\Ob{\mathop{\mathrm{Ob}}\nolimits}
\def\Sh{\mathop{\mathit{Sh}}\nolimits}
\def\NL{\mathop{N\!L}\nolimits}
\def\CH{\mathop{\mathrm{CH}}\nolimits}
\def\proetale{{pro\text{-}\acute{e}tale}}
\def\etale{{\acute{e}tale}}
\def\QCoh{\mathit{QCoh}}
\def\Ker{\mathop{\mathrm{Ker}}}
\def\Im{\mathop{\mathrm{Im}}}
\def\Coker{\mathop{\mathrm{Coker}}}
\def\Coim{\mathop{\mathrm{Coim}}}

% Boxtimes
%
\DeclareMathSymbol{\boxtimes}{\mathbin}{AMSa}{"02}

%
% Macros for moduli stacks/spaces
%
\def\QCohstack{\mathcal{QC}\!\mathit{oh}}
\def\Cohstack{\mathcal{C}\!\mathit{oh}}
\def\Spacesstack{\mathcal{S}\!\mathit{paces}}
\def\Quotfunctor{\mathrm{Quot}}
\def\Hilbfunctor{\mathrm{Hilb}}
\def\Curvesstack{\mathcal{C}\!\mathit{urves}}
\def\Polarizedstack{\mathcal{P}\!\mathit{olarized}}
\def\Complexesstack{\mathcal{C}\!\mathit{omplexes}}
% \Pic is the operator that assigns to X its picard group, usage \Pic(X)
% \Picardstack_{X/B} denotes the Picard stack of X over B
% \Picardfunctor_{X/B} denotes the Picard functor of X over B
\def\Pic{\mathop{\mathrm{Pic}}\nolimits}
\def\Picardstack{\mathcal{P}\!\mathit{ic}}
\def\Picardfunctor{\mathrm{Pic}}
\def\Deformationcategory{\mathcal{D}\!\mathit{ef}}

\setlength{\emergencystretch}{3em}
\begin{document}
\title[Stacks Project, Tag 0CRX]{380 --- Derived quasi-coherent pushforward agrees with sheaf-theoretic pushforward on affine-diagonal schemes}
\maketitle
\noindent Copyright (C) 2005--2025 Johan de Jong. Stacks Project authors. Source: \href{https://stacks.math.columbia.edu/tag/0CRX}{Tag 0CRX}.

\noindent Editorial difficulty note (not author text). Difficulty H2: The comparison is not a formal consequence of the underived equality because the two ambient categories have different resolutions. The proof identifies suitable pushed-forward K-injectives and then uses two kernel/image/cokernel sequences to induct on support and extend the comparison to arbitrary complexes..

\noindent Editorial provenance note (not author text). History: excerpt from revision \texttt{a0444\allowbreak 6e57e\allowbreak c1fbc\allowbreak 252a8\allowbreak 71afc\allowbreak ec775\allowbreak 2fb28\allowbreak 07b14}, perfect.tex, lines 1603--1694. Reference presentation was adapted; original-author context and core support blocks were retained or added. The mathematical wording is unchanged. Licensed under GNU FDL 1.2 or later; no invariant sections or cover texts. See ../licenses/COPYING. Context blocks reproduce author definitions and supporting results; they are not additional counted problems.

\begin{lemma}
\label{lemma-direct-image-coherator}
Let $f : X \to Y$ be a morphism of schemes.
Assume $X$ and $Y$ are quasi-compact and have affine diagonal.
Then, denoting
$$
\Phi : D(\QCoh(\mathcal{O}_X)) \to D(\QCoh(\mathcal{O}_Y))
$$
the right derived functor of
$f_* : \QCoh(\mathcal{O}_X) \to \QCoh(\mathcal{O}_Y)$
the diagram
$$
\xymatrix{
D(\QCoh(\mathcal{O}_X)) \ar[d]_\Phi \ar[r] &
D_\QCoh(\mathcal{O}_X) \ar[d]^{Rf_*} \\
D(\QCoh(\mathcal{O}_Y)) \ar[r] &
D_\QCoh(\mathcal{O}_Y)
}
$$
is commutative.
\end{lemma}

\begin{proof}
Observe that the horizontal arrows in the diagram are
equivalences of categories by
Proposition \href{https://stacks.math.columbia.edu/tag/08DB}{proposition quasi compact affine diagonal (Tag 08DB)}.
Hence we can identify these categories (and similarly for
other quasi-compact schemes with affine diagonal).
The statement of the lemma is that the canonical map
$\Phi(K) \to Rf_*(K)$ is an isomorphism for all $K$ in
$D(\QCoh(\mathcal{O}_X))$. Note that if $K_1 \to K_2 \to K_3 \to K_1[1]$
is a distinguished triangle in $D(\QCoh(\mathcal{O}_X))$ and
the statement is true for two-out-of-three, then it is true
for the third.

\medskip\noindent
Let $U \subset X$ be an affine open. Since the diagonal of $X$ is affine,
the inclusion morphism $j : U \to X$
is affine (Morphisms, Lemma \href{https://stacks.math.columbia.edu/tag/01SG}{lemma affine permanence (Tag 01SG)}).
Similarly, the composition $g = f \circ j : U \to Y$ is affine.
Let $\mathcal{I}^\bullet$ be a K-injective complex in $\QCoh(\mathcal{O}_U)$.
Since $j_* : \QCoh(\mathcal{O}_U) \to \QCoh(\mathcal{O}_X)$
has an exact left adjoint
$j^* : \QCoh(\mathcal{O}_X) \to \QCoh(\mathcal{O}_U)$
we see that $j_*\mathcal{I}^\bullet$ is a K-injective complex
in $\QCoh(\mathcal{O}_X)$, see
Derived Categories, Lemma \href{https://stacks.math.columbia.edu/tag/08BJ}{lemma adjoint preserve K injectives (Tag 08BJ)}.
It follows that
$$
\Phi(j_*\mathcal{I}^\bullet) =
f_*j_*\mathcal{I}^\bullet =
g_*\mathcal{I}^\bullet
$$
By Lemma \href{https://stacks.math.columbia.edu/tag/08D7}{lemma affine pushforward (Tag 08D7)} we see that
$j_*\mathcal{I}^\bullet$ represents $Rj_*\mathcal{I}^\bullet$ and
$g_*\mathcal{I}^\bullet$ represents $Rg_*\mathcal{I}^\bullet$.
On the other hand, we have $Rf_* \circ Rj_* = Rg_*$.
Hence $f_*j_*\mathcal{I}^\bullet$ represents $Rf_*(j_*\mathcal{I}^\bullet)$.
We conclude that the lemma is true for any complex
of the form $j_*\mathcal{G}^\bullet$ with $\mathcal{G}^\bullet$
a complex of quasi-coherent modules on $U$. (Note that if
$\mathcal{G}^\bullet \to \mathcal{I}^\bullet$ is a quasi-isomorphism,
then $j_*\mathcal{G}^\bullet \to j_*\mathcal{I}^\bullet$ is a
quasi-isomorphism as well since $j_*$ is an exact functor
on quasi-coherent modules.)

\medskip\noindent
Let $\mathcal{F}^\bullet$ be a complex of quasi-coherent
$\mathcal{O}_X$-modules. Let $T \subset X$ be a closed subset
such that the support of $\mathcal{F}^p$ is contained in $T$
for all $p$. We will use induction on the minimal number $n$
of affine opens $U_1, \ldots, U_n$ such that
$T \subset U_1 \cup \ldots \cup U_n$. The base case $n = 0$ is trivial.
If $n \geq 1$, then set $U = U_1$ and denote $j : U \to X$ the
open immersion as above. We consider the map of complexes
$c : \mathcal{F}^\bullet \to j_*j^*\mathcal{F}^\bullet$.
We obtain two short exact sequences of complexes:
$$
0 \to \Ker(c) \to \mathcal{F}^\bullet \to \Im(c) \to 0
$$
and
$$
0 \to \Im(c) \to j_*j^*\mathcal{F}^\bullet \to \Coker(c) \to 0
$$
The complexes $\Ker(c)$ and $\Coker(c)$ are supported
on $T \setminus U \subset U_2 \cup \ldots \cup U_n$ and the result
holds for them by induction. The result holds for
$j_*j^*\mathcal{F}^\bullet$ by the discussion in the preceding
paragraph. We conclude by looking at the distinguished triangles
associated to the short exact sequences and using the initial
remark of the proof.
\end{proof}
\end{document}
